\section{Limitations}
\label{sec:limitations}

\paragraph{One stream, one order, one model.}
Every number here comes from a single task order (Order-4) on a single backbone
(T5-large, LoRA $r{=}8$ on $q/v$) at a single budget ($400$/class, $3$ epochs).
Order effects in continual learning are large, and we have not measured them. The
scope structure that drives the whole analysis is a property of \emph{this}
benchmark: ten of fifteen tasks share output tokens. A stream whose tasks have
pairwise disjoint verbalizers has $\Delta^{\mathrm{id}}\equiv0$ by
Proposition~\ref{prop:disjoint}, and on such a stream nothing in this paper
applies. We regard the direction of that dependence as the honest framing: we
identify a term that exists when verbalizers are shared, and we say so rather than
implying generality.

\paragraph{The theory is proved in a special case that happens to cover us.}
Proposition~\ref{prop:qm} is proved for $d=1$ --- binary shared verbalizers --- and
for any $m$. The case $m\ge2$ with dimension $\ge2$ is \textbf{open}. It happens
that every scorable multi-task scope on this stream is binary, so every $q_m$ we
report is backed by proof; but that is a fact about which tasks survived our
exclusion rules, not a fact about the theory. A stream with scorable $K_S\ge3$
shared scopes would take us outside what we have proved.

\paragraph{The accuracy-space bound is weak where it matters.}
Proposition~\ref{prop:cheb-acc} carries a measured constant
$\kappa\approx0.018$ (worst task, WiC). At that value the bound is not
practically informative, and we do not rest any empirical claim on it: the
unconditional geometric statement does the work, and the accuracy numbers are
measured. Two conservatisms remain unresolved --- multi-class flip fractions count
only top-$2$ flips, and these $\kappa$ were estimated from an under-trained run
rather than re-estimated after convergence.

\paragraph{Three failed interventions are not a proof that none can work.}
Section~\ref{sec:three-anchors} reports that an orthogonality penalty of O-LoRA's
form, verbalizer-logit anchoring, and gauge-fixed anchoring all failed to improve
task-agnostic accuracy, two of them harmfully. The reading we prefer --- that the
identification gap is not a parameter-space object --- is an inference, not a
theorem. The alternative reading is that better parameter-space interventions exist
and we did not find them, and nothing in our evidence excludes it. We also tested
exactly one $\lambda$ grid per intervention, selected on a $6$-task prefix; in the
GFA case that gate demonstrably failed to predict $15$-task behaviour, which is a
weakness of our selection procedure and not only of the method.

\paragraph{Our two sensitivity analyses are underpowered, and we say so with numbers.}
Section~\ref{sec:sso} reports that neither the scale dependence of
$\Delta^{\mathrm{id}}$ nor the growth of staleness with stream length is established
here, and the design explains why. Live scope membership only ever takes the values
$\{1,2,3,4\}$, and only two scopes ever exceed one member, so task identity and
membership count cannot be fully separated; membership is also monotone in stage, so
``more scope-mates'' and ``later in the stream'' are collinear by construction. On
the staleness side, five scorable singleton tasks are trained at positions
$4,7,12,13,15$, giving age spans $0..11$, $0..8$, $0..3$, $0..2$ and $0..0$, so every
observation at age $\ge9$ comes from COPA alone. These are the reasons a failed
criterion here is weak evidence of absence rather than evidence of no effect, and
they were written down before the analysis was run.

\paragraph{Two of our four sensitivity arms are the same training trajectory.}
We had described Phase-2K and Phase-2P as two independent \texttt{none} arms for that
analysis. They are not: comparing their per-stage records field by field, all $480$
retention values agree \emph{bit for bit}, with identical gradient-step sequences and
identical seed, learning rate and epoch settings --- Phase-2P is Phase-2K's trajectory
with SSO scoring attached and its penalty weights at zero. The arm-invariance check is
therefore over one \texttt{none} replication (Phase-2Q) plus one different arm
(Phase-2M at $\lambda_a=0.5$), not three independent draws. The signs do agree across
all of them, but the check is weaker than we had claimed and we report it at the
corrected strength.

\paragraph{The enumeration is exhaustive only over its own constraint set.}
Section~\ref{sec:enumeration} claims exactly two non-vacuous admissible rules, and
that claim is an exhaustive case split over a table we wrote --- which examples a
rule may read and which model it may read them under. It is not a theorem about
task-agnostic inference. Three ways past it that our evidence does not close:
a rule that changes the training objective (our three anchors are three points in
that space, not a covering of it); a \emph{per-example} rather than per-task
router, which Section~\ref{sec:theory} shows can exceed even the per-task oracle
and whose task-free achievability we do not know; and stored state that is not an
estimate of a per-task optimum at all, for which our staleness mechanism gives no
prediction. ``Closed'' throughout this paper should be read as \emph{closed as
enumerated}, and the enumeration's premises are stated so they can be rejected.

\paragraph{Our O-LoRA arm is a re-implementation.}
The orthogonality penalty is our implementation of O-LoRA's regulariser inside our
harness at our budget. It is not the authors' code and not a reproduction of their
reported numbers. No comparison against their published table is made or implied,
and the arm exists to test additivity, not to benchmark.

\paragraph{The additivity question is unresolved, by our own criterion.}
Phase-2L's criterion C --- that the regulariser measurably repair deep forgetting on
the $>8$~pp stratum --- failed at both $\lambda$ values its gate selected. Our frozen
protocol says that when C fails, the additivity criteria A and B are
\emph{uninterpretable}. So we do not know whether an output-layer fix is additive
on top of a working representation-level method; we know only that it survived
alongside a penalty that did not measurably work. This is the single largest gap in
the paper's argument and we prefer to state it as such.

\paragraph{No comparison against published continual-PEFT methods.}
We did not run E$^2$-LoRA, NSR, LiteLoRA, SLICE, or ProCL. Every comparison here is
internal: arms differing in one component under an identical budget, identical
splits, and identical seeds. A claim that our approach outperforms published methods
would require running them, and we make no such claim.

\paragraph{Held-out data discipline costs us a number.}
The official test files are never opened. Everything is measured on \texttt{audit}
splits carved from training files, with offsets fit on disjoint \texttt{risk}
splits. This keeps the test sets clean for future work but means our absolute
accuracies are not comparable to published leaderboard numbers, and we do not place
them in such tables.

\paragraph{Our evaluation splits are class-balanced, and we measured what that buys
only once.}
Every \texttt{audit} split is $64$/class. Section~\ref{sec:bpo} shows this is not a
neutral choice: a rule that equalises predicted class frequencies gains $+4.40$~pp
on such a split and $-0.67$~pp at a $70{:}30$ mix. We pre-registered that control
for the query-batch rule because that rule's mechanism made the confound obvious.
\textbf{We did not run the equivalent control on anything else.} In particular the
$+1.51$~pp per-scope offset gain --- the paper's central measurement --- is a
balanced-split number, and while a per-scope offset fitted on labelled data has no
comparable reason to depend on query balance, we have not measured that it does not.
A reader should treat every accuracy in this paper as conditional on balanced
evaluation, and the one place we probed that condition, it mattered enough to change
a sign.

\paragraph{One frozen criterion could not be evaluated as written.}
Phase-2Q's protocol named a quantity its harness never recorded, so the criterion
guarding Proposition~\ref{prop:disjoint}'s wiring was unevaluable as frozen and we
reconstructed it --- correctly, and matching a formula already committed in an
earlier phase's decision script, but \emph{after} the phase's primary number was
known. We record it as unevaluable rather than passed, and the underlying prediction
is separately confirmed at $0/87$ by two other phases' checks. The general lesson is
against us: a pre-registration is only as strong as the harness that has to feed it,
and one of ours had a criterion that could never have fired. Our released notes
carry this and three other corrections to claims of our own, unedited beside the
appends that supersede them.

\paragraph{Excluded tasks.}
Three of fifteen tasks are excluded from scored means --- CB (rarest class $16$
examples, below the $40$ a $64$/class risk-plus-audit split needs), Yelp and Amazon
(five-way verbalizers sharing first sub-word pieces). The exclusions were fixed
before scoring and are the same in every arm, but they remove exactly the
higher-$K_S$ scopes, which is also where the open part of
Proposition~\ref{prop:qm} lives. A reader should treat our evidence as concerning
binary shared verbalizers.

\paragraph{Router state is not free, and we do not fold it into our headline cost.}
The $m{=}1$ table costs $33$ floats. Any $m>1$ variant additionally needs a
task-free router with per-task state that this figure does not cover; we report
that separately and never inside the same number.

\paragraph{Our proposed method is scored at two upper bounds.}
The combined method of Section~\ref{sec:combined} is measured with oracle scope
routing --- at evaluation we set each query's family adapter from its scope --- and
with the offset fit against the final model on each task's \texttt{risk} split. Both
are upper bounds, not deployable rules. A deployable system needs a task-free router
that recovers the family from the input alone; Section~\ref{sec:gap-real}'s honest
prototype router (nearest stored class-mean in encoder state) is the natural
candidate, and closing the grouped-oracle-to-grouped-routed gap is the first
measurement we owe. The $+11.76$~pp end-to-end number should therefore be read as the
achievable ceiling of the allocation, and a reviewer is entitled to discount it until
the router is measured. What is \emph{not} an upper bound is the fixed-budget
allocation itself: the four cells of the $2\times2$ share a bit-identical trainable
budget, so the comparison between them is exact regardless of routing.

\paragraph{Grouping de-stales the stored offset, but the deployable offset is not a net gain on its own.}
Section~\ref{sec:method} argued that a per-family adapter is perturbed only by its own
family, so a stored per-family offset should drift less than one stored against a
monolithic adapter. We measured this (Section~\ref{sec:combined}): grouping cuts the
staleness of the deployable stored offset from $8.22$~pp to $1.95$~pp, on all three
seeds, so the stored offset now tracks the refit upper bound the composed method reports.
Two honest limits remain. First, the offset column of the $2\times2$ is still \emph{run}
with the refit offset, not the stored one; we have shown the stored offset nearly matches
it, not re-run the whole $2\times2$ under the stored rule. Second, and more important, the
stored offset's net gain over no offset is $-0.51$~pp mean under grouping and mixed in
sign across seeds: grouping removes the active \emph{harm} that staleness caused on a
shared adapter ($-5.08$~pp in Section~\ref{sec:sso}), but on this stream it does not turn
the deployable offset into a reliable standalone improvement. The deployable value of the
offset is therefore that it composes with grouping at close to the upper bound, not that
it helps on its own.